\documentclass[coursework, och, times]{SCWorks1}
% параметр - тип обучения - одно из значений:
%    spec     - специальность
%    bachelor - бакалавриат (по умолчанию)
%    master   - магистратура
% параметр - форма обучения - одно из значений:
%    och   - очное (по умолчанию)
%    zaoch - заочное
% параметр - тип работы - одно из значений:
%    referat    - реферат
%    coursework - курсовая работа (по умолчанию)
%    diploma    - дипломная работа
%    pract      - отчет по практике
%    pract      - отчет о научно-исследовательской работе
%    autoref    - автореферат выпускной работы
%    assignment - задание на выпускную квалификационную работу
%    review     - отзыв руководителя
%    critique   - рецензия на выпускную работу
% параметр - включение шрифта
%    times    - включение шрифта Times New Roman (если установлен)
%               по умолчанию выключен
\usepackage{cmap}
\usepackage[T2A]{fontenc}
\usepackage[utf8]{inputenc}
\usepackage{graphicx}

\usepackage[sort,compress]{cite}
\usepackage{amsmath}
\usepackage{amssymb}
\usepackage{amsthm}
\usepackage{fancyvrb}
\usepackage{longtable}
\usepackage{array}
\usepackage[english,russian]{babel}


\usepackage[colorlinks=false]{hyperref}


\newcommand{\eqdef}{\stackrel {\rm def}{=}}

\newtheorem{lem}{Лемма}
\newtheorem{theorem}{Теорема}

\begin{document}

% Кафедра (в родительном падеже)
\chair{теории функции и стохастического анализа}

% Тема работы
\title{Сходимость несобственных интегралов: признаки, методы исследования и приложения}

% Курс
\course{2}

% Группа
\group{212}

% Факультет (в родительном падеже) (по умолчанию "факультета КНиИТ")
\department{механико-математического факультета}

% Специальность/направление код - наименование
\napravlenie{01.03.02 "--- Прикладная математика и информатика}
%\napravlenie{02.03.02 "--- Фундаментальная информатика и информационные технологии}
%\napravlenie{02.03.01 "--- Математическое обеспечение и администрирование информационных систем}
%\napravlenie{09.03.01 "--- Информатика и вычислительная техника}
%\napravlenie{09.03.04 "--- Программная инженерия}
%\napravlenie{10.05.01 "--- Компьютерная безопасность}

% Для студентки. Для работы студента следующая команда не нужна.
%\studenttitle{Студентки}

% Фамилия, имя, отчество в родительном падеже
\author{Чесакова Максима Евгеньевича}

% Заведующий кафедрой
\chtitle{д.\,ф.-м.\,н., доцент} % степень, звание
\chname{С.\,П.\,Сидоров}

%Научный руководитель (для реферата преподаватель проверяющий работу)
\satitle{доцент, д.\,ф.-м.\,н.} %должность, степень, звание
\saname{С.\,П.\,Сидоров}


% Год выполнения отчета
\date{2026}

\maketitle
% Включение нумерации рисунков, формул и таблиц по разделам
% (по умолчанию - нумерация сквозная)
% (допускается оба вида нумерации)
%\secNumbering
\tableofcontents

% Раздел "Обозначения и сокращения". Может отсутствовать в работе
%\abbreviations
%\begin{description}
%    \item $|A|$  "--- количество элементов в конечном множестве $A$;
%    \item $\det B$  "--- определитель матрицы $B$;
%    \item ИНС "--- Искусственная нейронная сеть;
%    \item FANN "--- Feedforward Artifitial Neural Network
%\end{description}

% Раздел "Определения". Может отсутствовать в работе
%\definitions

% Раздел "Определения, обозначения и сокращения". Может отсутствовать в работе.
% Если присутствует, то заменяет собой разделы "Обозначения и сокращения" и "Определения"
%\defabbr


% Раздел "Введение"
\intro
введение

\section{Основные понятия и определения}

\subsection{Несобственные интегралы I рода}
\textbf{Определение 1}. \textit{Пусть функция $f(x)$ определена на $[a,\infty)$ и $\forall b>a$ $f(x)$ интегрируема на $[a,b]$. Тогда величина
\[
\int\limits_a^\infty f(x)\,dx = \lim_{b \to \infty}\int\limits_a^b f(x)\,dx
\]
называется несобственным интегралом Римана I рода.
}

Если указанный предел существует, то интеграл называется сходящимся и значение предела --- это значение интеграла. Иначе интеграл называется расходящимся. Таким образом, вопрос о сходимости несобственного интеграла равносилен вопросу о том, определен ли вообще этот несобственный интеграл или нет.

\subsection{Несобственные интегралы II рода}
\textbf{Определение 2}. \textit{
Пусть функция $f(x)$ определена на $[a,b)$ и \\ $\displaystyle\lim_{x \to b -0} f(x) = \infty$, $\forall b'\in[a,b)$ $f(x)$ интегрируема на $[a,b']$. Тогда величина
\[
\int\limits_a^b f(x)\,dx = \lim_{b' \to b-0}\int\limits_a^{b'} f(x)\,dx
\]
называется несобственным интегралом Римана II рода.
}

Суть этого определения состоит в том, что в любой окрестности конечной точки $b$ функция $f$ может оказаться неограниченной.

\subsection{Абсолютная и условная сходимость}


\section{Признаки сходимости несобственных интегралов}

\subsection{Интегральный признак Коши}


\subsection{Признаки сравнения}


\subsection{Признак Дирихле и признак Абеля}


\subsection{Критерий Коши сходимости несобственного интеграла}


\subsection{Эквивалентные преобразования подынтегральной функции}


\section{Методы вычисления несобственных интегралов}

\subsection{Интегрирование по частям и замена переменной в несобственных интегралах}


\subsection{Вычисление через пределы: сведение к собственному интегралу}


\subsection{Использование специальных функций}


\subsection{Метод дифференцирования по параметру}


\section{Примеры исследования и вычисления несобственных интегралов}

\subsection{Интеграл Пуассона $\int_{0}^{\infty} e^{-x^{2}} \, dx$ и его вычисление}




%==================================================================================================================================================================================
%||||||||||||||||||||||||||||||||||||||||||||||||||||||||||||||||||||||||||||||||||||||||||||||||||||||||||||||||||||||||||||||||||||||||||||||||||||||||||||||||||||||||||||||||||
%==================================================================================================================================================================================
%||||||||||||||||||||||||||||||||||||||||||||||||||||||||||||||||||||||||||||||||||||||||||||||||||||||||||||||||||||||||||||||||||||||||||||||||||||||||||||||||||||||||||||||||||
%==================================================================================================================================================================================
%----------------------------------------------------------------------------------------------------------------------------------------------------------------------------------
%==================================================================================================================================================================================
%||||||||||||||||||||||||||||||||||||||||||||||||||||||||||||||||||||||||||||||||||||||||||||||||||||||||||||||||||||||||||||||||||||||||||||||||||||||||||||||||||||||||||||||||||
%==================================================================================================================================================================================
%||||||||||||||||||||||||||||||||||||||||||||||||||||||||||||||||||||||||||||||||||||||||||||||||||||||||||||||||||||||||||||||||||||||||||||||||||||||||||||||||||||||||||||||||||
%==================================================================================================================================================================================


% Раздел "Заключение"
\conclusion
В ходе выполнения курсовой работы были исследованы аналитические методы теории телетрафика и их практическое применение для анализа систем массового обслуживания. Основное внимание уделено математическому аппарату теории вероятностей и теории марковских процессов, которые составляют фундамент теории телетрафика.

Ключевым достижением работы стала программная реализация аналитической модели системы M/M/1 на языке Python. Разработанный программный комплекс позволяет вычислять стационарные вероятности и основные характеристики системы, включая среднее число заявок в системе и очереди, среднее время пребывания и ожидания. Визуализация зависимостей этих характеристик от параметров $\lambda$ и $\mu$ наглядно демонстрирует нелинейный рост времени ожидания и длины очереди при приближении коэффициента загрузки к единице.

Практическая часть подтвердила эффективность аналитических методов для решения задач теории телетрафика в случаях относительно простых структур систем. Полученные результаты полностью соответствуют теоретическим предсказаниям и наглядно иллюстрируют критические режимы работы систем массового обслуживания.

Несмотря на ограниченную применимость аналитических методов для сложных систем, их сочетание с численными методами и статистическим моделированием остается перспективным направлением для анализа телекоммуникационных систем. Важно подчеркнуть, что математические модели всегда приближенно описывают реальные системы, поэтому полученные результаты следует рассматривать как теоретическую основу для практического проектирования.

Проведенное исследование демонстрирует практическую ценность аналитических методов теории телетрафика и возможность их успешной программной реализации для анализа характеристик систем массового обслуживания.
%Библиографический список, составленный вручную, без использования BibTeX
%
%\begin{thebibliography}{99}
%  \bibitem{Ione} Источник 1.
%  \bibitem{Itwo} Источник 2
%\end{thebibliography}

\begin{thebibliography}{99}

  \bibitem{1} Лившиц\,Б.\,С. Теория телетрафика : учебник для вузов / Б.\,С.\,Лившиц, А.\,П.\,Пшеничников, А.\,Д.\,Харкевич. -- 2-е изд., перераб. и доп. -- М.~: Связь, 1979. -- 224 с.
  \bibitem{2} Зорин\,А.\,В. Введение в общие цепи Маркова : Учебно-методическое пособие / В.\,А.\,Зорин, Е.\,В.\,Пройдакова, М.\,А.\,Федоткин -- Нижний Новгород~: Нижегородский госуниверситет, 2013. — 51 с.
  \bibitem{3} Крылов\,В.\,В. Теория телетрафика и ее приложения / В.\,В.\,Крылов, С.\,С.\,Самохвалова. -- СПб.~: БХВ-Петербург, 2005. -- 288 с.
  \bibitem{4} Кельберт\,М.\,Я. Вероятность и статистика в примерах и задачах. Т.II: Марковские цепи как отправная точка теории случайных процессов и их приложения. / Ю.\,М.\,Сухов -- М.~: МЦНМО, 2009. — 588 с.
  \bibitem{5} Наместников\,С.\,М. Основы теории телетрафика : учебное пособие / С.\,М.\,Наместников, М.\,Н.\,Служивый, Ю.\,Д.\,Украинцев. -- Ульяновск~: УлГТУ, 2016. -- 154 с.
  \bibitem{6} Бочаров\,П.\,П. Теория вероятностей. Математическая статистика / П.\,П.\,Бочаров, А.\,В.\,Печенкин. -- M.~: Наука, Физматгиз, 2005. -- 296~с.
  \bibitem{7} Гмурман\,В.\,Е. Теория вероятностей и математическая статистика : Учебное пособие для вузов. / В.\,Е.\,Гмурман. -- М.~: Высш. шк., 2003. -- 479~с.
  \bibitem{8} Гладков\,Л.\,Л. Теория вероятностей и математическая статистика : Учебное пособие / Л.\,Л.\,Гладков, Г.\,А.\,Гладкова. -- Мн.~: РИПО, 2013. -- 248~с.
  \bibitem{9} Python [Электронный ресурс] : язык программирования Python : официальная документация, версия 3.x / Python Software Foundation. -- Электрон. дан. -- [США], 2001-. -- Режим доступа: https://docs.python.org/3, свободный. -- Загл. с экрана.
  \bibitem{10} Буйначев\,С.\,К. Основы программирования на языке Python : учебное пособие / С.\,К.\,Буйначев, Н.\,Ю.\,Боклаг. – Екатеринбург~: Изд-во Урал. ун-та, 2014. -- 91 c. 

\end{thebibliography}


%Библиографический список, составленный с помощью BibTeX
%
%\bibliographystyle{gost780uv}
%\bibliography{thesis}

% Окончание основного документа и начало приложений
% Каждая последующая секция документа будет являться приложением
\appendix

\section{Листинг программы}\label{pril}
\begin{Verbatim}[fontsize=\small, numbers=left]
import numpy as np
import matplotlib.pyplot as plt

# Настройка для отображения графиков и формул
plt.rcParams['font.size'] = 12
plt.rcParams['figure.figsize'] = (10, 6)

def mm1_characteristics(lambd, mu):
    if lambd >= mu:
        print(f"Предупреждение: lambda ({lambd}) >= mu ({mu})."
              f"Система нестабильна (rho >= 1).")
        return None
    
    rho = lambd / mu  # Коэффициент загрузки
    
    P0 = 1 - rho            # Вероятность простоя системы
    N = rho / (1 - rho)     # Среднее число заявок в системе
    Nq = N * rho            # Среднее число заявок в очереди
    T = 1 / (mu - lambd)    # Среднее время в системе
    W = T * rho             # Среднее время в очереди
    
    # Вероятности для первых 10 состояний
    k_states = 10
    Pk = [(1 - rho) * (rho ** n) for n in range(k_states + 1)]
    
    characteristics = {
        'rho': rho,
        'P0': P0,
        'N': N,
        'Nq': Nq,
        'T': T,
        'W': W,
        'Pk': Pk
    }
    return characteristics

# 1. Анализ зависимости характеристик от lambda (при фиксированном mu)
mu_fixed = 1.0  # Фиксируем интенсивность обслуживания
# lambda меняется от почти 0 до почти mu
lambda_values = np.linspace(0.01, 0.99, 50)

# Массивы для хранения результатов
N_values = []
Nq_values = []
T_values = []
W_values = []
rho_values = []

for lambd in lambda_values:
    chars = mm1_characteristics(lambd, mu_fixed)
    if chars is not None:
        rho_values.append(chars['rho'])
        N_values.append(chars['N'])
        Nq_values.append(chars['Nq'])
        T_values.append(chars['T'])
        W_values.append(chars['W'])

# Построение графиков для зависимости от lambda
fig, axes = plt.subplots(2, 2, figsize=(14, 10))
fig.suptitle('Зависимость характеристик СМО M/M/1 '
             'от интенсивности lambda (mu = 1.0 фикс.)', fontsize=16)

# График 1: Среднее число заявок
axes[0, 0].plot(lambda_values, N_values, 'b-',
                linewidth=2, label='N (в системе)')
axes[0, 0].plot(lambda_values, Nq_values, 'r--',
                linewidth=2, label='Nq (в очереди)')
axes[0, 0].set_xlabel('Интенсивность lambda')
axes[0, 0].set_ylabel('Среднее число заявок')
axes[0, 0].set_title('Среднее число заявок')
axes[0, 0].legend()
axes[0, 0].grid(True, alpha=0.3)

# График 2: Среднее время
axes[0, 1].plot(lambda_values, T_values, 'g-',
                linewidth=2, label='T (в системе)')
axes[0, 1].plot(lambda_values, W_values, 'm--',
                linewidth=2, label='W (в очереди)')
axes[0, 1].set_xlabel('Интенсивность lambda')
axes[0, 1].set_ylabel('Среднее время')
axes[0, 1].set_title('Среднее время пребывания')
axes[0, 1].legend()
axes[0, 1].grid(True, alpha=0.3)

# График 3: Коэффициент загрузки
axes[1, 0].plot(lambda_values, rho_values, 'k-', linewidth=2)
axes[1, 0].set_xlabel('Интенсивность lambda')
axes[1, 0].set_ylabel('rho')
axes[1, 0].set_title('Коэффициент загрузки rho')
axes[1, 0].grid(True, alpha=0.3)

# График 4: Вероятность простоя системы
P0_values = [1 - rho for rho in rho_values]
axes[1, 1].plot(lambda_values, P0_values, 'c-', linewidth=2)
axes[1, 1].set_xlabel('Интенсивность lambda')
axes[1, 1].set_ylabel('P0')
axes[1, 1].set_title('Вероятность простоя системы P0')
axes[1, 1].grid(True, alpha=0.3)

plt.tight_layout(rect=[0, 0, 1, 0.96])
plt.show()

# 2. Анализ зависимости характеристик от mu (при фиксированном lambda)
lambda_fixed = 0.7  # Фиксируем интенсивность входящего потока
# mu должно быть больше lambda. Возьмем диапазон от lambda+0.1 до 3.0
mu_values = np.linspace(lambda_fixed + 0.1, 3.0, 50)

# Массивы для хранения результатов
N_values_mu = []
T_values_mu = []
Nq_values_mu = []
W_values_mu = []

for mu in mu_values:
    chars = mm1_characteristics(lambda_fixed, mu)
    if chars is not None:
        N_values_mu.append(chars['N'])
        Nq_values_mu.append(chars['Nq'])
        T_values_mu.append(chars['T'])
        W_values_mu.append(chars['W'])

# Построение графиков для зависимости от mu
fig, (ax1, ax2) = plt.subplots(1, 2, figsize=(14, 5))
fig.suptitle('Зависимость характеристик СМО M/M/1 от '
             'интенсивности обслуживания mu (lambda = 0.7 фикс.)',
             fontsize=16)

# График 1: Среднее число заявок от mu
ax1.plot(mu_values, N_values_mu, 'b-',
        linewidth=2, label='N (в системе)')
ax1.plot(mu_values, Nq_values_mu, 'r--',
        linewidth=2, label='Nq (в очереди)')
ax1.set_xlabel('Интенсивность обслуживания mu')
ax1.set_ylabel('Среднее число заявок')
ax1.set_title('Среднее число заявок')
ax1.legend()
ax1.grid(True, alpha=0.3)

# График 2: Среднее время в системе от mu
ax2.plot(mu_values, T_values_mu, 'g-',
         linewidth=2, label='T (в системе)')
ax2.plot(mu_values, W_values_mu, 'm--',
         linewidth=2, label='W (в очереди)')
ax2.set_xlabel('Интенсивность обслуживания mu')
ax2.set_ylabel('Среднее время пребывания')
ax2.set_title('Среднее время пребывания')
ax2.legend()
ax2.grid(True, alpha=0.3)

plt.tight_layout(rect=[0, 0, 1, 0.96])
plt.show()

# 3. Визуализация стационарного распределения вероятностей
# Рассчитаем распределение вероятностей для конкретных lambda и mu
lambd_example = 0.8
mu_example = 1.0
chars_example = mm1_characteristics(lambd_example, mu_example)

if chars_example is not None:
    k_states = len(chars_example['Pk'])
    states = np.arange(k_states)

    # График распределения вероятностей Pk
    plt.figure(figsize=(10, 6))
    plt.plot(states, chars_example['Pk'], 'bo', linewidth=2, markersize=6)
    plt.xlabel('Число заявок в системе, k')
    plt.ylabel('Вероятность, Pk')
    plt.title(f'Стационарное распределение вероятностей состояний СМО M/M/1\n'
              f'(lambda={lambd_example}, mu={mu_example})')
    plt.grid(True, alpha=0.3)
    plt.xticks(states)
    plt.show()

    # Вывод численных характеристик для примера
    print("\nХарактеристики системы для lambda={} и mu={}:".format(
        lambd_example, mu_example))
    print("Коэффициент загрузки rho: {:.4f}".
          format(chars_example['rho']))
    print("Среднее число заявок в системе N: {:.4f}".
          format(chars_example['N']))
    print("Среднее время пребывания в системе T: {:.4f}".
          format(chars_example['T']))
    q = 0
    for i in chars_example['Pk']:
        print(q, ' & ', end='')
        print("{:.6f}".format(i))
        q+=1
    print("summ Pk = {:.6f}".format(sum(chars_example['Pk'])))
\end{Verbatim}


\end{document}
